\documentclass[../DoAn.tex]{subfiles}
\begin{document}

Phụ lục này mô tả cấu hình và nội dung tổng quát của các prompt dùng cho module so khớp CV--JD và module phân tích CV trong hệ thống. Nội dung đầy đủ của prompt cùng schema JSON được công bố trong mã nguồn của khóa luận.

\section{Cấu hình LLM dùng cho module so khớp CV--JD}

Bảng~\ref{tab:llm-config-ch5} liệt kê các tham số gọi mô hình \texttt{gpt-4o-mini} của OpenAI thông qua giao thức Chat Completions trong module so khớp.

\begin{table}[H]
\centering
\caption{Tham số gọi LLM trong module so khớp CV--JD}
\label{tab:llm-config-ch5}
\begin{tabular}{|l|l|}
\hline
\textbf{Tham số} & \textbf{Giá trị} \\
\hline
Model & \texttt{gpt-4o-mini} \\
\hline
max\_tokens & 400 \\
\hline
temperature & 0 \\
\hline
seed & 42 \\
\hline
response\_format & \texttt{json\_object} \\
\hline
Provider & OpenAI Chat Completions API \\
\hline
Ngôn ngữ prompt & tiếng Anh (khớp với dataset Kaggle) \\
\hline
\end{tabular}
\end{table}

\subsection{Nội dung prompt}

Prompt gồm hai phần: system prompt thiết lập vai trò và thang điểm, và user prompt chứa cặp CV–JD cần đánh giá.

System prompt yêu cầu mô hình đóng vai một chuyên viên tuyển dụng kỹ thuật có kinh nghiệm, chấm điểm độ phù hợp giữa CV và JD trên thang 0–100 với bốn mức diễn giải rõ ràng:

\begin{itemize}
\item 90--100: phù hợp tốt -- hầu hết kỹ năng yêu cầu và kinh nghiệm liên quan đều có.
\item 70--89: phù hợp khá -- có đủ kỹ năng cốt lõi nhưng còn vài thiếu sót.
\item 40--69: phù hợp một phần -- có một số điểm trùng nhưng cũng có nhiều khoảng trống quan trọng.
\item 0--39: ít phù hợp -- khác lĩnh vực hoặc thiếu kinh nghiệm đáng kể.
\end{itemize}

System prompt đồng thời yêu cầu mô hình chỉ trả về JSON hợp lệ, không kèm văn xuôi.

User prompt nối JD và CV theo cấu trúc cố định, đóng khung bằng nhãn (\texttt{=== JOB DESCRIPTION ===} và \texttt{=== CANDIDATE CV ===}) để mô hình phân biệt hai đầu vào. Phần cuối yêu cầu trả về một đối tượng JSON gồm bốn trường: điểm tổng \texttt{overallScore}, danh sách từ khóa trùng \texttt{matchedKeywords}, danh sách từ khóa thiếu \texttt{missingKeywords}, và một câu giải thích ngắn \texttt{reasoning}.

\subsection{Schema đầu ra}

Bảng~\ref{tab:llm-schema} mô tả các trường trong JSON đầu ra của LLM. Module so khớp parse JSON, ép kiểu và giới hạn miền giá trị cho từng trường trước khi trả về kết quả.

\begin{table}[H]
\centering
\caption{Schema JSON đầu ra của module so khớp CV--JD}
\label{tab:llm-schema}
\begin{tabular}{|l|l|p{7cm}|}
\hline
\textbf{Trường} & \textbf{Kiểu} & \textbf{Mô tả} \\
\hline
\texttt{overallScore} & integer [0,100] & Điểm tổng độ phù hợp 0--100. \\
\hline
\texttt{matchedKeywords} & mảng chuỗi (tối đa 10) & Các kỹ năng/từ khóa xuất hiện ở cả CV và JD. \\
\hline
\texttt{missingKeywords} & mảng chuỗi (tối đa 10) & Từ khóa quan trọng trong JD nhưng thiếu ở CV. \\
\hline
\texttt{reasoning} & chuỗi (1 câu) & Lý do ngắn gọn cho điểm số. \\
\hline
\end{tabular}
\end{table}

Khi mô hình trả về văn bản không đúng định dạng (ví dụ kèm theo dấu \texttt{```json...```}), wrapper sẽ tiền xử lý loại bỏ phần thừa rồi parse lại. Nếu vẫn không hợp lệ, cặp đánh giá được đánh dấu lỗi và ghi vào cột \texttt{error} của tệp kết quả \texttt{results.csv}.

\subsection{Lưu ý về tính tái lập}

Tham số \texttt{seed = 42} kết hợp \texttt{temperature = 0} giúp giảm tối đa biến thiên giữa các lần chạy. Tuy nhiên, OpenAI ghi nhận seed chỉ ở mức ``best-effort'' do phần cứng GPU và cơ chế cân tải phía máy chủ; mức biến thiên còn lại được trường \texttt{system\_fingerprint} trong phản hồi của API ghi nhận lại. Khóa luận ghi nhận yếu tố này như một hạn chế trong phần đánh giá tính tin cậy của kết quả.

\section{Cấu hình prompt trong module phân tích CV của hệ thống}

Module \texttt{jd-analyze} của hệ thống ngoài việc chấm điểm còn sinh gợi ý cải thiện cho từng phần của CV nên có prompt và schema phong phú hơn so với mục~B.1. Bảng~\ref{tab:llm-compare} liệt kê các điểm khác biệt chính.

\begin{table}[H]
\centering
\caption{So sánh cấu hình LLM giữa module so khớp và module phân tích JD}
\label{tab:llm-compare}
\begin{tabular}{|l|l|l|}
\hline
\textbf{Tham số} & \textbf{Module so khớp} & \textbf{Module jd-analyze} \\
\hline
max\_tokens & 400 & 3000 \\
\hline
temperature & 0 & 0.3 \\
\hline
seed & 42 & (không đặt) \\
\hline
Ngôn ngữ & tiếng Anh & song ngữ Việt -- Anh (theo locale) \\
\hline
Đầu ra mong đợi & chỉ điểm số + từ khóa & điểm số + từ khóa + gợi ý chi tiết \\
\hline
\end{tabular}
\end{table}

System prompt của module sản phẩm yêu cầu mô hình đóng vai một chuyên gia HR Việt Nam có 15 năm kinh nghiệm. Trong tiếng Việt, system prompt nhấn mạnh yêu cầu phản hồi bằng tiếng Việt và trả về JSON hợp lệ; trong tiếng Anh, prompt giữ vai trò ``experienced HR professional'' tương tự nhưng không ràng buộc về ngôn ngữ phản hồi.

\subsection{Schema gợi ý cải thiện}

Bảng~\ref{tab:jd-analyze-schema} mô tả các nhóm trường chính của JSON đầu ra. Mỗi gợi ý là một bản ghi (record) gồm tiêu đề, mô tả, mức độ ưu tiên (high/medium/low), nội dung đề xuất, và các từ khóa liên quan.

\begin{table}[H]
\centering
\caption{Cấu trúc đầu ra của module jd-analyze}
\label{tab:jd-analyze-schema}
\begin{tabular}{|l|p{9cm}|}
\hline
\textbf{Nhóm trường} & \textbf{Mô tả} \\
\hline
\texttt{analysisId} & Định danh duy nhất của phân tích. \\
\hline
\texttt{jobMatch} & Điểm tổng, từ khóa trùng/thiếu, điểm mạnh và điểm cần cải thiện. \\
\hline
\texttt{suggestions} & Bốn nhóm gợi ý theo từng phần CV: summary, experience, skills, education; tối đa 3 gợi ý mỗi nhóm. \\
\hline
\texttt{globalRecommendations} & Danh sách khuyến nghị tổng thể không gắn vào phần nào cụ thể. \\
\hline
\end{tabular}
\end{table}

Vì payload đầu ra dài (\texttt{max\_tokens = 3000}), prompt của module phân tích JD tạo ra lượng token lớn hơn đáng kể so với các module khác trong hệ thống.

\section{Các prompt khác trong hệ thống}

Bảng~\ref{tab:other-llm} liệt kê các điểm API còn lại có sử dụng LLM trong hệ thống.

\begin{table}[H]
\centering
\caption{Các điểm API khác sử dụng LLM trong hệ thống}
\label{tab:other-llm}
\begin{tabular}{|l|p{7cm}|c|}
\hline
\textbf{Điểm API} & \textbf{Mục đích} & \textbf{max\_tokens} \\
\hline
\texttt{rewrite-bullet} & Viết lại 1 bullet trong phần kinh nghiệm. & 200 \\
\hline
\texttt{wizard-bullet} & Sinh bullet từ mô tả công việc tự do của người dùng. & 300 \\
\hline
\texttt{summary} & Sinh đoạn tóm tắt chuyên môn từ phần còn lại của CV. & 400 \\
\hline
\texttt{score-feedback} & Chấm điểm CV theo bốn tiêu chí: ATS, nội dung, cấu trúc, ngôn ngữ. & 600 \\
\hline
\end{tabular}
\end{table}

Các prompt này được thiết kế ngắn nhằm tiết kiệm chi phí khi vận hành sản phẩm. Nội dung đầy đủ của từng prompt cùng schema đầu ra có thể tham khảo trong mã nguồn của hệ thống.

\end{document}
